% Transcription — fonds Alexandre Grothendieck, shelfmark 15, batch 3 (pages 41-60).
% Established from the facsimile archives/batches/15/batch-03.pdf (a mirror of
% https://grothendieck.umontpellier.fr/15.pdf), per the skill
% transcribe-grothendieck. The batch file carries no cover sheet: PDF page N
% is page 40+N of the folder (the Montpellier footer prints « 41 » ... « 60 »),
% and the archivists' pencil numbers bottom-left agree on every transcribed
% leaf (41, 43, 45, 47, 49, 51, 53, 55, 58, 59, 60 all legible).
%
% Pass: Opus 5.5 (claude-opus-5-5), 2026-09-23 — first pass, unchecked against the pages by a human.
% (A retry: a first pass on this batch stalled before writing; its scratch
% draft of pages 45-55 was consulted and every page re-read against the
% tiles.)
%
% Eleven of the twenty pages are transcribed. Absent: pages 42, 44, 46, 48,
% 50, 52, 54 and 56, the versos of pages 41-55 — leaves of the unrelated
% typescript on holomorphic functions of several variables already met in
% batches 2 and 4 (typed pp. 9-24, several upside down), nothing in his hand;
% page 57, blank.
%
% What the batch is. Two pieces.
%   (1) Pages 41-55, yellowed paper, handwritten, not paginated by him,
%   continuing batch 2's Deuring categories of supersingular elliptic curves
%   (the quaternion algebra over Q ramified at p and ∞). Page 41 does not
%   visibly continue the « On trouve donc » that ends page 39.
%   41     the split Deuring category (C'(A'), T') over A' = ∏_{ℓ≠p,∞} Z_ℓ
%          (square roots of A'(1)), C(Z_p) = F-crystals of rank 2, slope 1/2
%          over F̄_p, C(R) via the gerbe g_∞: a torsor T'(C) under
%          Γ' = A'*/A'*² × (Z_p*/Z_p*² × Z/2Z) × R*/R*²; the Frobenius
%          autoequivalence acts by (p, f = (0,1), p); T(C) = T'(C)/±1, a
%          torsor under Γ ≃ Γ'_0 = H¹(adèles, μ_2).
%   43     T(F̄_p) a canonical torsor under H¹(A, μ_2) = ∏_ℓ Z_ℓ*/Z_ℓ*² (ℓ
%          including p and ∞); Ẑ acts trivially, so the torsor is canonical.
%          NB: suggests a trivial μ_2-gerbe on A with T as its torsor of
%          classes; pairs (C, α), α a class of equivalences
%          (C,T)_{Z_p} ≃ (C(Z_p), T(Z_p)) modulo Pic(Z_p). Two margin notes
%          (the two isomorphism classes (M,o), (M^(p),o) swapped by Frobenius;
%          « Trop compliqué » — α as an o ⊗ Z_p-structure on a crystal).
%   45     any two (C,α) are equivalent, the equivalences forming a connected
%          groupoid with π_1 = Z/2Z; the gerbe G(C,α) over A = ∏ Z_ℓ, of
%          group μ_2, determined up to a unique-up-to-sign equivalence but
%          depending on F̄_p; F̄_p = lim F_p[T]/Φ_{p^n−1}(T).
%   47     the 2-category S of the (C,α) with π_0 = π_1 = 0, π_2 = Z/2Z; the
%          chain Éq(C_π) → Éq(C,T) → Éq(C_p,T_p) ≃ Éq(C(Z_p),T(Z_p)), fully
%          faithful; the classes of equivalences modulo Éq(C) form a torsor
%          T_p under Z_p*/Z_p*²; ξ ∈ T_p.
%   49     a canonical gerbe G(ξ) on the adèles (on A'' = A' × A_∞ the
%          equivalences of G_{Z[1/p]} with g'_{A''}, trivial at p); the theory
%          of elliptic curves of Hasse invariant 0 gives a) the « élément de
%          Deuring » ξ_p ∈ T_p, b) a section ξ̂ of G(ξ_p), up to sign at ∞.
%   51     « À l'infini, la situation n'est pas très belle encore »: the gerbe
%          h(ξ_p)_R, trivial, polarised, « mais il reste l'ambiguïté de
%          signe »; C(ξ_p)_{/∞} ≃ one-dimensional H-vector spaces (H the
%          quaternions); a functor « canonique au signe près » to R.
%   53     « Détermination de ξ_p ? »: an element ξ of T_p as a function
%          (M, o) ↦ u_ξ(M,o) ∈ Isom(Λ²M, T_p)/Z_p*² (M an F-crystal of rank 2
%          over F_{p²} with o-action, T_p the Tate crystal), conditions
%          (i) stability under isomorphism, (ii) compatibility with twisting
%          by an invertible P (det M' ≃ N_{o/Z}(P) ⊗ det M ≃ det M),
%          a triangle diagram; NB: (i) amounts to Frobenius-invariance.
%   55     equivalently triples (M_0, o, i), i : R → o, R the integers of
%          Q[T]/(T² + p), i.e. a descent of M to F_p with F² + p = 0.
%   (2) Pages 58-60, white paper, headed « Catégories de Deuring. », the
%   folder's own typescript:
%   58     typed: « 1. Catégories A-linéaires et ⊗-enveloppes. »: an
%          A-linear functor is « admissible » (compatible with external
%          tensor product) if P ⊗_o F(X) → F(P ⊗_o X) is an isomorphism;
%          for a Deuring category test it on one X with o = End_C(X) and
%          P ∈ P(o/A); Hom_{A-adm}(C, C') ≃ objects of C' with o-action;
%          C(1) = Hom_{A-adm}(C°, Mod(A)) and the Yoneda embedding
%          i : C → C(1). His annotations: a typed marginal « notion duale
%          "coadmissible" », a brown-ink note on a « générateur strict »,
%          and three handwritten lines at the foot.
%   59     entirely in his hand, continuing page 58: C_sp(1), functors whose
%          values are projective of finite type (« spéciaux »), i : C →
%          C_sp(1); a strict generator X identifies C_sp(1) with projective
%          right o-modules of finite type, the « saturée » of C.
%   60     in his hand: (1) the map P ⊗_o X → Hom_o(P̌, X); external tensor
%          products « stricts », « raisonnables »; P ↦ P ⊗_o X fully
%          faithful; (2) « catégorie jolie ». The page ends mid-sentence,
%          « Si C admet un gén. strict X, (tout objet », continued on page 62
%          (batch 4).
%
% Notation fixed for the batch. o for the ring of endomorphisms (as in batch
% 4; batch 2 writes σ for the same letter); 𝒪 the set of such orders;
% \mathbb{A} for the adèles (his A with a bar), \mathbb{A}', \mathbb{A}'' as
% he writes them (on page 41 he writes A' for ∏_{ℓ≠p,∞} Z_ℓ); μ_2 for the
% « fl2 »-like glyph (the group of the gerbes); T for the Tate object,
% 𝒯, 𝒯', 𝒯_p for the torsors, 𝒢 for the gerbes, 𝔥(ξ_p) for his gothic h;
% Éq as in batch 2 for his underlined sigle; \mathbb{P}(o/A) as in batch 2.
% Plain square brackets in the text are his: either his own brackets, or
% words he added in the interline; editorial additions are \add{}.
%
% Legibility. The formulas carry; the fast connective prose and the struck
% passages of pages 43-51 are partly \ill{}; page 51 is in a paler ink and
% partly in pencil. Pages 53-55 and 58-60 read almost end to end. Nothing on
% these leaves is dated.
%
% The dating is the inventory's own for the folder, copied verbatim. Its
% group, [10-18] « Théorie des motifs », is dated [à partir de 1958-à partir
% de 1982].
\documentclass[11pt,a4paper]{article}
% Le préambule commun à toutes les transcriptions du fonds Grothendieck.
%
% Il tient en une idée : **l'incertitude se compose, elle ne se résout pas.**
% Une transcription qui lisse un mot douteux a détruit la seule chose pour
% laquelle on la faisait — un lecteur doit pouvoir savoir, à chaque endroit,
% ce qui était sur le feuillet et ce qui a été ajouté. Les six macros qui
% suivent sont donc l'appareil critique, et non de la décoration.
%
% Les mêmes commandes sont reconnues par `scripts/render.mjs`, qui produit la
% vue de lecture du site. Une macro ajoutée ici doit l'être aussi là-bas,
% faute de quoi le rendu échouera bruyamment — ce qui est voulu : un rendu
% silencieusement faux est pire qu'un rendu absent.

\ProvidesPackage{grothendieck}[2026/08/08 Transcriptions du fonds Grothendieck]

\usepackage{amsmath,amssymb,amsthm}
\usepackage{mathrsfs}
\usepackage{stmaryrd}
% Les diagrammes commutatifs sont partout dans ce fonds : c'est la langue
% même de la géométrie algébrique de Grothendieck, et une transcription qui
% les rendrait en prose ne transcrirait rien.
\usepackage{tikz-cd}
% Français par défaut — c'est la langue des feuillets — mais l'anglais est
% chargé aussi : la traduction et le résumé sont des documents anglais, et la
% typographie française (espaces avant les deux-points) y serait une faute.
% Ces documents basculent par \selectlanguage{english} en tête de corps.
\usepackage[english,french]{babel}
\usepackage{fontspec}
\usepackage{geometry}
\geometry{a4paper,margin=25mm,marginparwidth=28mm}
\usepackage{marginnote}
\usepackage{xcolor}
% ulem, not soul. `soul`'s \st and \ul are built on a character-by-character
% scan that XeLaTeX's font machinery breaks: the document fails with an
% undefined control sequence rather than degrading. ulem does the same two jobs
% and is Unicode-engine safe. `normalem` stops it hijacking \emph, which the
% transcriptions use for Grothendieck's own emphasis.
\usepackage[normalem]{ulem}
\usepackage{hyperref}
\hypersetup{colorlinks=true,linkcolor=black,urlcolor=black,pdfborder={0 0 0}}

% --- Métadonnées du lot -----------------------------------------------------
% Reprises telles quelles par le script de rendu, qui les affiche en tête de
% la vue de lecture. Elles fixent ce que le fichier prétend couvrir : sans
% elles, une transcription flotte sans point d'attache dans un fonds de seize
% mille feuillets.

\newcommand{\folder}[1]{\gdef\grothFolder{#1}}
\newcommand{\batch}[1]{\gdef\grothBatch{#1}}
\newcommand{\leaves}[2]{\gdef\grothFirst{#1}\gdef\grothLast{#2}}
\newcommand{\foldertitle}[1]{\gdef\grothFoldertitle{#1}}
\gdef\grothFolder{?}\gdef\grothBatch{?}\gdef\grothFirst{?}\gdef\grothLast{?}\gdef\grothFoldertitle{}

% --- Le résumé --------------------------------------------------------------
%
% Il ouvre la lecture modernisée, sous le titre « Résumé », et il est composé
% en petit. Ce n'est pas un ornement : c'est le seul passage du document qui
% ne soit pas des mathématiques, et un lecteur doit pouvoir le reconnaître
% avant de l'avoir lu — puis le sauter s'il connaît déjà le sujet, ou s'y
% arrêter s'il ne le connaît pas. Le filet qui le suit marque la reprise du
% corps.

\newenvironment{resume}
  {\small\setlength{\parskip}{0.45em}}
  {\par\medskip\hrule\medskip}

% --- Mots-clés ---------------------------------------------------------------
%
% En anglais, en clôture du résumé de la lecture modernisée : le vocabulaire
% moderne sous lequel chercher ce que les feuillets construisent. Le manifeste
% du site (`npm run manifest`) les extrait pour étiqueter la cote entière —
% cette ligne est la source unique des tags, il n'y a pas de fichier à côté.
\newcommand{\keywords}[1]{\par\smallskip\noindent
  {\footnotesize\textsc{Keywords} --- \textit{#1}}\par}

% --- L'appareil critique ----------------------------------------------------

% Le numéro de feuillet, en marge. C'est l'ancre qui permet de régler un
% désaccord sur une formule en regardant *un* feuillet plutôt qu'une chemise
% de six cents. Le site s'en sert aussi pour tourner les pages du fac-similé
% quand on fait défiler la transcription.
\newcommand{\leaf}[1]{%
  \marginnote{\footnotesize\color{gray!70}\textsf{#1}}%
  \label{leaf:#1}\ignorespaces}

% Illisible : on ne devine pas. Un mot inventé qui se lit comme les autres est
% le pire résultat possible.
\newcommand{\ill}{\textcolor{red!60!black}{[\,\ldots\,]}}

% Lecture douteuse : le mot est donné, et signalé comme incertain.
\newcommand{\uncertain}[1]{\uline{#1}}

% Ajout de l'éditeur — une lettre manquante, un mot sous-entendu.
\newcommand{\add}[1]{\textcolor{blue!50!black}{[#1]}}

% Biffé par Grothendieck lui-même. Ce qu'il a rayé fait partie du document :
% une rature dit souvent où la pensée a changé de direction.
\newcommand{\struck}[1]{\sout{#1}}

% Note du transcripteur, sur l'état matériel du feuillet.
\newcommand{\note}[1]{\par\smallskip{\small\color{gray!60!black}\textit{[#1]}}\par\smallskip}

% Note marginale de Grothendieck, distincte du corps du texte.
\newcommand{\marginal}[1]{\par\smallskip\noindent\hspace{1em}%
  {\small\color{gray!60!black}#1}\par\smallskip}

% --- Titre ------------------------------------------------------------------

\AtBeginDocument{%
  \begin{center}
    {\large\bfseries Fonds Alexandre Grothendieck --- cote n\textsuperscript{o}~\grothFolder}\\[2pt]
    {\itshape\grothFoldertitle}\\[4pt]
    {\small feuillets \grothFirst--\grothLast\ (lot \grothBatch)}\\[2pt]
    {\footnotesize\color{gray!60!black}Transcription établie d'après le fac-similé numérisé
     par l'Université de Montpellier.\\
     Les lectures incertaines sont soulignées, les illisibles notées [\,\ldots\,],
     les ajouts entre crochets.}
  \end{center}
  \vspace{1em}}

\endinput


\folder{15}
\batch{3}
\pages{41}{60}
\dating{1965-[vers 1977]}
\watermark{Édition de démonstration}
\foldertitle{Théorie arithmétique. Théorie de Galois des motifs (1965) : notes manuscites (s.d.), tapuscrit (s.d.).}

\begin{document}

\page{41}
Ici on a donné la \struck{\ill{}} $(C'(A'), T')$
[catégorie de Deuring \uncertain{déployée}] sur
$A' = \prod_{\ell \neq p, \infty} \mathbb{Z}_{\ell}$ comme associée à la
gerbe des racines carrées de $A'(1)$ (définie par le choix d'une clôture
algébrique $\overline{\mathbb{F}}_{p}$ de $\mathbb{F}_{p}$), et la
catégorie de Deuring $C(\mathbb{Z}_{p})$ comme celle des $F$-cristaux sur
$\overline{\mathbb{F}}_{p}$ [avec le $F$-cristal de Tate]
\uncertain{(à iso.)} qui sont de rang~$2$, \struck{\ill{}} pente
$\tfrac{1}{2}$, et la \struck{gerbe} catégorie
$C(\mathbb{Z}_{\infty}) = C(\mathbb{R})$ comme associée à la gerbe
$g''_{\infty}$ \ill{} sur $\mathbb{R}$, on trouve un torseur
\struck{\ill{}} [$\mathcal{T}'(C)$] \uncertain{certain} torseur sous le
groupe
\[
  A'^{*}/A'^{*2} \times \bigl(\mathbb{Z}_{p}^{*}/\mathbb{Z}_{p}^{*2}
  \times \mathbb{Z}/2\mathbb{Z}\bigr) \times \mathbb{R}^{*}/\mathbb{R}^{*2}
  = \Gamma' .
\]
\note{un premier $\mathbb{Z}_{p}$, biffé, précède la formule. La lecture de l'indice de $g''_{\infty}$ et des mots qui suivent est
incertaine}

Le groupe [$\mathbb{Z}/2\mathbb{Z}$] des \struck{classes} d'isomorphie
d'autoéquivalences de \struck{$C$} \ill{} opère sur ce torseur, de façon
compatible avec son opération triviale sur \ill{}, donc par translations,
via l'image $f$ [(Frobenius)] de $\mathbb{Z}/2\mathbb{Z}$ dans ce groupe,
et il opère donc via \struck{\ill{}} l'image [\uncertain{gén.}] $f$ de
$\mathbb{Z}/2\mathbb{Z}$ dans les $3$ facteurs [\uncertain{respectivement}]
\[
  \bigl(p, \; f = (0, 1), \; p\bigr)
\]
\marginal{N.B. \struck{On peut dire}, [comme $p \in \mathbb{R}^{*2}$]
\struck{plus \ill{}}, \ill{} l'un des \ill{},
\struck{distingué canonique de} \struck{$\mathrm{Pl}(C)$} il y a un système
transitif d'isom.\ plusieurs~!}
\note{la note N.B. est rattachée par une flèche au troisième terme $p$ ; ses
deux premières lignes sont raturées en plusieurs couches}

Considérons le quotient
\[
  \Gamma = \Gamma' / f(\pm 1)
\]
\struck{\ill{}} et soit $\mathcal{T}(C)$ le quotient de $\mathcal{T}'(C)$
par $\pm 1$, qui est donc un torseur sous $\Gamma$. Ce torseur [à isom.~
unique près] ne dépend pas de $C$, [i.e.\ \uncertain{on identifie}
$\mathcal{T}$ à $\mathcal{T}(\overline{\mathbb{F}}_{p})$]. D'ailleurs si
$\Gamma'_{0}$ est le sous-groupe de $\Gamma'$
\[
  A'^{*}/A'^{*2} \times \mathbb{Z}_{p}^{*}/\mathbb{Z}_{p}^{*2} \times
  \mathbb{R}^{*}/\mathbb{R}^{*2} = \mathrm{H}^{1}(\mathbb{A}, \mu_{2})
\]
\note{sous $\mathbb{A}$, en petit : « adèles » ; $\mathbb{A}$ surcharge un
premier signe biffé}
on trouve que $\Gamma' = \Gamma'_{0} \times f(\pm 1)$, donc
$\Gamma \simeq \Gamma'_{0}$ can., donc
\note{la phrase se poursuit au feuillet 43}

\page{43}
\[
  \mathcal{T}(\overline{\mathbb{F}}_{p}) \ \text{est un torseur canonique
  sous} \ \mathrm{H}^{1}(\mathbb{A}, \mu_{2}) =
  \prod_{\ell} \mathbb{Z}_{\ell}^{*}/\mathbb{Z}_{\ell}^{*2} \quad
  (\text{y compris } p \text{ et } \infty) .
\]
\note{le signe $=$ est écrit verticalement sous $\mathrm{H}^{1}(\mathbb{A},
\mu_{2})$, le produit au-dessous ; « (y compris $p$ et $\infty$) » est en
petit sous l'indice $\ell$}

Ce torseur \uncertain{ici} dépend encore de $\overline{\mathbb{F}}_{p}$,
\struck{\ill{}} et cette dépendance est précisée par la connaissance des
opérations de $\widehat{\mathbb{Z}} \ni f$ sur
$\mathcal{T}(\overline{\mathbb{F}}_{p})$. Cette opération se fait via
[les opérations habituelles sur $A'(1)$, \struck{\ill{}}] les opérations de
$\widehat{\mathbb{Z}}$ sur $C(A', T')$, \struck{[\ill{}]}
[$C(A_{p}, T_{p})$] (via opération de \struck{$\mathbb{Z}$} $f$ sur
$C(A_{p}, T_{p})$ par Frobenius) et l'opération triviale sur
$C(A_{\infty}, T_{\infty})$. Donc on trouve que $f \in \widehat{\mathbb{Z}}$
opère par le $=$ élément $(p, f, p)$ que plus haut, donc
\[
  f \ \text{opère trivialement sur} \ \mathcal{T}(\overline{\mathbb{F}}_{p}) .
\]
Ainsi, on a un torseur canonique [$\mathcal{T}$] sous
$\mathrm{H}^{1}(\mathbb{A}, \mu_{2})$.

\marginal{\ill{} : l'objet : quaternions $(M, o)$ : \emph{isom.} \ill{} il y
a deux classes d'isomorphie, \ill{} $(M, o)$ est l'une d'elles, l'autre est
$(M^{(p)}, o)$ \ill{} \uncertain{distinguées} par Frobenius}
\note{note marginale à gauche, en haut du feuillet, écrite le long du bord
et encadrée ; un trait la relie au NB qui suit. Le premier mot est
surchargé}

[\,[ \emph{NB} Ceci suggère qu'il y a une \struck{\ill{}, canonique}
[gerbe triviale (un] \struck{trivialisée}) de \struck{faisceau}
$\mu_{2}$ sur $\mathbb{A}$, dont $\mathcal{T}$ serait le torseur des
classes d'isomorphie. On peut \uncertain{Plutôt} considérons les couples
$(C, \struck{\ill{}}\,\alpha)$, [d'où $T_{C} = T$,] où $C$ est une
catégorie de Deuring de type $\xi \in \operatorname{Br}(\mathbb{Q})$, et
$\alpha$ une \struck{\ill{}} classe d'équivalences des classes d'isomorphie
[d'équivalences] de $(C, T)_{\mathbb{Z}_{p}}$ avec
$(C(\mathbb{Z}_{p}), T(\mathbb{Z}_{p}))$, modulo les équivalences provenant
de $\mathrm{Pic}(\mathbb{Z}_{p})$ ; i.e.\ une
section du torseur de groupe $\mathbb{Z}/2\mathbb{Z}$ déduit de celui de
groupe
$\text{Éq}(C(\mathbb{Z}_{p}), T(\mathbb{Z}_{p})) =
\mathbb{Z}_{p}^{*}/\mathbb{Z}_{p}^{*2} \times \mathbb{Z}/2\mathbb{Z}$ par
extension du groupe structural à l'aide de $\mathrm{pr}_{2}$.
\note{le sigle devant $(C(\mathbb{Z}_{p}), T(\mathbb{Z}_{p}))$ est lu
« Éq » par analogie avec le lot 2 ; la lecture est incertaine.
$\mathrm{Pic}$ est souligné deux fois}

\marginal{Trop compliqué : il suffit de fixer une classe d'isomorphie
d'équivalences de cat.\ de Deuring $C_{\mathbb{Z}_{p}} \simeq
C(\mathbb{Z}_{p})$ --- en termes concrets, un $\alpha$ est donné (si $C$
défini par $o$) en se donnant une \uncertain{structure} de
$o \otimes_{\mathbb{Z}} \mathbb{Z}_{p}$ sur l'anneau des \uncertain{endom.}~
d'un \ill{} cristal sur $\overline{\mathbb{F}}_{p}$ (de rang $2$,
\uncertain{pente} $\tfrac{1}{2}$ \ill{}), \struck{\ill{}} \ill{} de Tate}
\note{note marginale à gauche, en bas du feuillet, écrite en oblique ; les
deux dernières lignes sont raturées et entourées}

Alors \struck{\ill{}} $C$ admet un $\alpha$, et deux $(C, \alpha)$,
$(C', \alpha')$
\note{la phrase se poursuit au feuillet 45}

\page{45}
\struck{sont \uncertain{canoniquement} équivalentes, [l'équivalence donnée
à isom.\ près] Si on a [\ill{}] un couple essentiellement identifiés entre
eux les différents $(C, \alpha)$, et on a ainsi une catégorie de Deuring
\emph{canonique})}
\note{les trois premières lignes et le mot « canonique », souligné, sont
encadrés et biffés de traits obliques}
sont équivalents : la catégorie des équivalences de $(C, \alpha)$ avec
$(C', \alpha')$ est un groupoïde connexe de \struck{$\pi_{1}$,}
$\mathbb{Z}/2\mathbb{Z}$ : i.e.\ \struck{\ill{}} toutes les équivalences
sont isom., et si on a deux équiv.\ $\varphi, \varphi'$, il y a \emph{à
signe près} un seul isom.\ entre $\varphi$ et $\varphi'$.

\struck{Soit donné un couple $(C, \alpha)$}

Soit donné maintenant un couple $(C, \alpha)$, il lui est associé une
\struck{torseur sous} gerbe [de $(C, T)$] sur
$\mathbb{A} = \prod_{\ell} \mathbb{Z}_{\ell}$ [$\ell$ y compris $\ell = p$
et $\ell = \infty$] de groupe $\mu_{2}$, [savoir] celle des équivalences
avec le \struck{gerbe} couple type $CC(\mathbb{A}, T)$, qui en la place $p$
soit compatible avec la donnée de $\alpha$. Cette gerbe
$\mathcal{G}(C, \alpha)$ est déterminée, à une « signe près » \struck{\ill{}}
équivalence unique près, indépendamment du choix de $(C, \alpha)$
(puisqu'il en est ainsi des $(C, \alpha)$ lui-même $= \cdots$),
\struck{et $\mathcal{G} = \mathrm{Gal}(\overline{\mathbb{F}}_{p}/
\mathbb{F}_{p})$ y opère \uncertain{trivialement}.} mais elle dépend
[$=$ pour $(C, \alpha)$ fixé, du choix] d'un $\overline{\mathbb{F}}_{p}$
\struck{(ou du moins d'un $\mathbb{F}_{p^{2}}$)} --- tout comme la
structure elle-même~!
\note{« $CC(\mathbb{A}, T)$ » est lu tel quel}

\marginal{\struck{Car le $F$-cristal de rang $2$, pente $\tfrac{1}{2}$ sur
$\overline{\mathbb{F}}_{p}$ se descend \uncertain{canoniquement} à
$\mathbb{F}_{p^{2}}$}}
\note{note marginale à gauche, biffée de traits obliques}

On veut s'en tirer en définissant \emph{canoniquement} un
$\mathbb{F}_{p^{2}}$ (ou même un $\overline{\mathbb{F}}_{p}$, si on y
tenait …) par les polynômes cyclotomiques $\Phi_{p^{n}-1}$,
\[
  \overline{\mathbb{F}}_{p} = \varinjlim_{n}
  \bigl( \mathbb{F}_{p^{n}} = \mathbb{F}_{p}[T]/\Phi_{p^{n}-1}(T) \bigr) \qquad
  \text{-----}
\]

\page{47}
Considérons cette gerbe $\mathcal{G}$ \struck{sur} [sur]
$\mathbb{A} = \prod_{p} \mathbb{Z}_{p}$, de groupe $\mu_{2}$, définie « au
signe près » (de façon précise, on a une $2$-catégorie $\mathcal{S}$ des
$(C, \alpha)$, qui satisfait $\pi_{0} = \pi_{1} = 0$,
$\pi_{2} = \mathbb{Z}/2\mathbb{Z}$, et \struck{\ill{}} on a [une famille de
gerbes] \struck{\ill{}} indexée par $\mathcal{S}$, i.e.\ un $2$-foncteur de
$\mathcal{S}$ dans la $2$-catégorie des \struck{$2$-}gerbes sur
$\mathbb{A}$, groupe $\mu_{2}$). Dans cette gerbe, on a une « section
définie à isom.\ can.\ au signe près ».\,]\,]
\note{les deux dernières lignes sont soulignées et fermées par un double
crochet dans la marge de droite}
\marginal{attention : \uncertain{ici bis}, \ill{} un deuxième signe~!}

\note{un court trait de séparation}

Autre façon d'expliquer : la \struck{catégorie des Éq} [foncteur des]
catégories de Picard \uncertain{comp.}
\[
  \text{Éq}(C_{\pi}) \longrightarrow \text{Éq}(C, T) \longrightarrow
  \text{Éq}(C_{p}, T_{p}) \underset{\mathrm{can}}{\simeq}
  \text{Éq}\bigl(C(\mathbb{Z}_{p}), T(\mathbb{Z}_{p})\bigr)
\]
\note{le sigle, souligné, est lu « Éq » comme au feuillet 43 ; l'indice de
$C_{\pi}$ est douteux ; le dernier terme est écrit sous
$\text{Éq}(C_{p}, T_{p})$, relié par « $\simeq$ can »}
est pleinement fidèle, la catégorie de Picard quotient est
\struck{\ill{}} une catégorie rigide, qui n'est autre que
$\mathcal{Z}[\,0 \to \mathbb{Z}_{p}^{*}/\mathbb{Z}_{p}^{*2}\,]$. Donc la
catégorie des équivalences de $(C, T)_{\mathbb{Z}_{p}}$ avec
$\bigl(C(\mathbb{Z}_{p}), T(\mathbb{Z}_{p})\bigr)$ $\longrightarrow$
\emph{modulo action de} $\text{Éq}(C)$ est une catégorie rigide,
\struck{torseur} dont l'ens.\ des classes d'isom.\ est un torseur
\struck{canonique} sous $\mathbb{Z}_{p}^{*}/\mathbb{Z}_{p}^{*2}$ (on devrait
\struck{la gerbe} [dire un vrai $\mathcal{T}_{p}$] !!!), indépendant à
isom.\ can.\ près de $C$. \struck{On en déduit \ill{} pour \ill{} à can.}
$\xi \in \mathcal{T}_{p}$.
\struck{\ill{} pour considérer \ill{} d'une \ill{}} les
\struck{Grâce à ce choix,} \ill{} $\alpha$ est une équivalence
$(C, \alpha)$, i.e.\ \ill{}
$(C, T)_{\mathbb{Z}_{p}} \simeq \bigl(C(\mathbb{Z}_{p}), T(\mathbb{Z}_{p})\bigr)$
\emph{dans la classe de} $\xi_{p}$, on \ill{} dans telle gerbe, elles sont
\struck{isomorphes} équiv.\ \emph{de façon canonique}.
\note{la fin du feuillet est raturée en plusieurs couches, et l'ordre des
mots y est incertain ; « modulo action de » est souligné. Le symbole
$\mathcal{Z}$ devant le crochet est douteux}

\page{49}
Donc on peut \struck{considérer que} [identifier] les $(C, \alpha)$
\struck{\ill{}} entre eux \struck{\ill{}}, en tenant compte des
\[
  \bigl[\,C(\xi), \alpha(\xi)\,\bigr] \quad \text{canoniques.}
\]
\struck{On peut donc considérer sur $\mathbb{A}''$
[$= \mathbb{A}' \times \mathbb{A}_{\infty}$] (anneau des adèles premiers à
$p$), [de groupe $\mu_{2}$] la gerbe $g''$ qui sur $\mathbb{A}'$ est la
gerbe des racines carrées de $A'(1)$, et sur $\mathbb{A}_{\infty}$ est la
gerbe type $g'_{\infty}$ déjà envisagée,}
\note{passage encadré et biffé de traits obliques}

En fonction de ce choix, on trouve une gerbe \emph{canonique}
$\mathcal{G}(\xi)$ sur l'anneau des adèles, de groupe $\mu_{2}$, savoir la
gerbe qui sur $\mathbb{A}'' = \mathbb{A}' \times \mathbb{A}_{\infty}$ est
la gerbe des équivalences de la gerbe $\mathcal{G}_{\mathbb{Z}[1/p]}$
[de groupe $\mu_{2}$] définie sur $\mathbb{Z}[\tfrac{1}{p}]$ par $(C, T)$
\struck{\ill{}} avec la gerbe $g'_{\mathbb{A}''}$ déjà envisagée plus haut
(sur $\mathbb{A}'$ c'est la $\sqrt{\phantom{x}}$ de $A'(1)$, sur
$\mathbb{A}_{\infty}$ c'est $g'_{\infty}$ déjà envisagée plus haut) et qui
sur $\mathbb{A}_{p}$ est la gerbe \struck{dite} [triviale].
\struck{[sections ; isom.] Cette gerbe est la triviale, et l'une des
\ill{} $\mathcal{T}$ sous $\mathrm{H}^{1}(\mathbb{A}, \mu_{2})$ dit,}
\struck{Cette gerbe \ill{} près \ill{} équiv.\ \ill{} définie par \ill{}
$\mathcal{G}(\xi)$ de $\mathcal{G}(\xi)$} \struck{section canonique}
[de Deuring] [envisagée plus haut, qui est indép.\ de $\xi$.]
\note{le bas de ce passage est encadré et raturé en plusieurs couches ; on
a gardé ce qui se lit. Dans la marge de gauche, deux mots biffés, \ill{}}

Tout ceci \uncertain{posé}, la théorie des [courbes elliptiques]
\uncertain{d'inv.\ de Hasse nul} nous fournit
\begin{itemize}
  \item[a)] un élément $\xi_{p} \in \mathcal{T}_{p}$ (l'\emph{élément de
  Deuring})
  \item[b)] une section \struck{correspondante} $\hat{\xi}$ de
  $\mathcal{G}(\xi_{p})$, définie [\emph{au signe près à l'infini}]
  \struck{\ill{}} à isom.\ près, unique \ill{}, dont la valeur en $p$
  \struck{\ill{}} est la section unité, et la valeur en $\infty$ est la
  [polarisation donnée par $T$ ….] \struck{\ill{}} de
  $\mathcal{G}(\xi_{p})_{\mathbb{Z}_{p}}$, [et] $|\xi|$ de $\mathcal{T}$,
  pour la
\end{itemize}
On trouve donc un [élément] \struck{section} [dont la] compos.\ suivant $p$
est $\xi_{p}$.
\note{« de Deuring » est écrit au-dessus de « la théorie des », sans qu'on
voie à quoi il se rattache ; la fin de b) est raturée et se lit mal}

\page{51}
À l'infini, la situation n'est pas très belle encore. Je sais seulement
dire :

Pour tout $\xi_{p} \in \mathcal{T}_{p}$, [considérons la gerbe
$\mathfrak{h}(\xi_{p})$] \struck{la valeur en $\infty$ de
$\mathcal{G}(\xi)$} de groupe $\mu_{2}$ sur $\mathbb{Z}_{1/p}$
définie par $\bigl(C(\xi_{p}), T(\xi_{p})\bigr)$, \struck{\ill{}}
[considérons] $\mathfrak{h}(\xi_{p})_{\mathbb{R}}$, qu'elle est une gerbe
\emph{canonique} sur $\mathbb{R}$ de groupe $\mu_{2}$ [pour $\xi_{p} =$
élément de Deuring,] \struck{j'ignore \ill{}} [elle a une description
indépendante] de la théorie des catégories de Deuring (et de celle des
courbes elliptiques). Ceux sont en tous cas \ill{} qu'elle est triviale ;
elle est, pour $\xi_{p} =$ élément de Deuring, munie d'une polarisation
$\xi_{\infty}$ (ce qui revient à \ill{} d'une classe d'isomorphie
[d'équivalences] avec la gerbe type $g_{\infty}$), \emph{mais il reste
l'ambiguïté de signe}.
\note{page écrite d'une encre plus pâle, au crayon par endroits ; plusieurs
mots de liaison se perdent. « $\mathbb{Z}_{1/p}$ » est écrit ainsi, pour
$\mathbb{Z}[\tfrac{1}{p}]$}

Catégorie $C(\xi_{p})_{/\infty}$ \struck{---} $\approx$ catégorie des
[droites $V$ de dim.~$1$] vectoriels \struck{à gauche} sur $\mathbb{H}$
($=$ quaternions ordinaires), avec \struck{\ill{}} $V$ foncteur norme
ordinaire vers cat.\ de Picard de $\mathbb{R}$.

Ainsi il y a un foncteur « canonique au signe près »
\[
  C(\xi_{p}) \xrightarrow{\ \varphi_{C_{\infty}}\ }
  \mathcal{C}\ell(\mathbb{R})
\]
\note{le but surcharge un premier but biffé, \struck{$C(\mathbb{H})$
catégorie} ; sous lui, $\mathbb{H}$ ; la lecture
$\mathcal{C}\ell(\mathbb{R})$ du but et celle de l'indice de la flèche sont
incertaines}
\marginal{\uncertain{muni par ailleurs} de $C(\xi_{p})$ pour tout $C$}
\note{note à gauche de la formule, reliée à elle par une flèche}
\marginal{(le foncteur \ill{} dans les vectoriels quaternioniens de dim.~$1$
\emph{unitaires} a deux isom.\ \uncertain{chiales}~!)}
\struck{$(T(\xi_{p})$ \ill{} défini \ill{} n'est autre} \ill{} [à un autre
choix] défini \emph{au signe près}, \ill{}, l'isom.

Prenant $\xi_{p} =$ élément de \struck{Frobenius} Deuring, on trouve un
foncteur [ess.\ canonique] \struck{des} [courbes elliptiques] de H.\ nul
[\uncertain{droite unitaire}] vers les vectoriels de dim.~$1$. La
\struck{\ill{}} Norme, la transposition correspond : à polarisation, la
transposition
\note{l'écriture se défait dans les deux dernières lignes ; la fin du
feuillet est blanche. « H.\ nul » est lu « de Hasse nul », comme au
feuillet 49}

\page{53}
\emph{Détermination de $\xi_{p}$ ?}

Explicitons d'abord $\mathcal{T}_{p}$. Un élément [$\xi$] de
$\mathcal{T}_{p}$ est une fonction qui à tout couple $(M, o)$, d'un
$o \in \mathcal{O}$ et d'un $F$-cristal $M$ sur $\mathbb{F}_{p^{2}}$ (ou
même sur $\overline{\mathbb{F}}_{p}$) de rang~$2$, muni d'une opération de
$o$, associe un \struck{isom} élément de
\[
  u_{\xi}(M, o) \in
  \mathrm{Isom}\bigl(\textstyle\bigwedge^{2} M, T_{p}\bigr) /
  \mathbb{Z}_{p}^{*2}
\]
\note{sous $T_{p}$, en petit : « cristal de Tate »}
de façon à ce que les conditions suivantes soient satisfaites :
\begin{itemize}
  \item[(i)] Stabilité par isomorphismes $(M, o) \to (M', o')$
  \item[(ii)] Pour $(M, o)$ et $P$ un $o$-module à droite projectif de
  rg~$1$, si $o'$ [$= \mathrm{End}_{o}(P)$] est son commutant, on a le
  couple
  \[
    (\underbrace{P \otimes_{o} M}_{M'}, o')
  \]
  et on a alors un isom.\ can.\ de $F$-cristaux
  \[
    \det M' \simeq (\mathrm{N}_{o/\mathbb{Z}} P) \otimes_{\mathbb{Z}}
    \det M \underset{c}{\simeq} \det M
  \]
  grâce à l'isom.\ can.\ $\mathrm{N}_{o/\mathbb{Z}}(P) \simeq \mathbb{Z}$.
  On veut que l'\ill{} \uncertain{commute} \ill{}
\end{itemize}

\begin{tikzcd}[column sep=small]
  \det(M) \arrow[rr, "c"] \arrow[dr, no head, "u_{\xi}(M{,}o)"'] & &
  \det M' \arrow[dl, no head, "u_{\xi}(M'{,}o')"] \\
  & T_{p} &
\end{tikzcd}

\note{les deux côtés obliques sont des traits sans pointe ; celui de droite
porte « $\approx$ », celui de gauche un signe qui ressemble à un « $2$ » en
exposant, peut-être le même « $\approx$ » ; la flèche horizontale porte
« $c$ » au-dessus de « $\simeq$ »}

\emph{NB} En présence de (ii), \struck{(i)} (i) interprété au sens strict
[\struck{$\otimes$ désignant n'importe} (en termes de $o, o' \in
\mathcal{O}$, bimodules ${}_{o}P_{o'}$ …, et [donné] isom.\ entre
$({}_{o}P_{o'}) \otimes_{\mathbb{Z}} \mathbb{Z}_{p}$ et le bimodule
$\mathrm{Hom}_{\mathrm{cris}}(M, M')$] la condition (i) équivaut à la
condition de stabilité par isom.\ $\bar{k} \to \bar{k}'$, ou encore par
l'automorphisme de Frobenius de $\bar{k}$ (ou de $\mathbb{F}_{p^{2}}$) ….

\page{55}
Il est équivalent de se donner [les] $u_{\xi}$ pour un triple
$(M_{0}, o, i)$, où \struck{$(M, o)$ \ill{}} $M_{0}$ est un cristal
\struck{sur} $\mathbb{F}_{p^{2}}$, $o$ \struck{\ill{} opérant} opérant sur
[$M = M_{0} \otimes \mathbb{F}_{p^{2}}$], et $i : R \to o$ une inclusion
($R$ l'anneau des entiers de $\mathbb{Q}[T]/(T^{2} + p)$) ; il revient au
même de se donner $i$, ou une descente [$M_{0}$] de $M$ à $\mathbb{F}_{p}$,
telle que [Frobenius sur $M_{0}$ satisfasse $F^{2} + p = 0$]
\struck{telle que $R$ opère sur $M_{0}$} (i.e.\ son opération sur $M$ se
descend à $\mathbb{F}_{p}$). Les conditions sont
\begin{itemize}
  \item[(i)] stabilité par isom.\ de triples
  $(M_{0}, o, i) \simeq (M'_{0}, o', i')$
  \item[(ii)] \struck{Stab} compatibilité avec twist par un module
  inversible $P_{0}$ sous $R$, compte tenu de l'isom.\ can.~
  $\det_{\mathbb{Z}} P_{0} \simeq \mathbb{Z}$.
\end{itemize}
\note{l'expression entre crochets $M = M_{0} \otimes \mathbb{F}_{p^{2}}$
surcharge une première version raturée ; la moitié inférieure du feuillet
est blanche}

\section{Catégories de Deuring}

\page{58}
\note{feuillet dactylographié, qui fait partie du tapuscrit du dossier ;
les ajouts à la main de l'auteur sont donnés entre crochets ou en note
marginale, et signalés. Le titre ci-dessus est celui du tapuscrit}

\subsection{1. Catégories $A$-linéaires et $\otimes$-enveloppes.}

$A$ est un anneau commutatif. Soit $C$ une catégorie $A$-linéaire (i.e.~
les Hom dans $A$ sont des $A$-modules, l'accouplement des Hom étant
bilinéaire). Un foncteur $A$-linéaire
\[
  F : C \longrightarrow C'
\]
de $C$ dans une autre catégorie $A$-linéaire $C'$ est dit « admissible » ou
« compatible au produit tensoriel externe » si pour tout objet $X$ de $C$
sur lequel opère ($A$-linéairement) une $A$-algèbre $o$, et tout
$o$-module à droite $P$ tel que $P \otimes_{o} X$ existe dans $C$,
$P \otimes_{o} F(X)$ existe dans $C'$ et l'homomorphisme naturel
\[
  P \otimes_{o} F(X) \longrightarrow F(P \otimes_{o} X)
\]
est un isomorphisme.
\note{dans la marge de gauche, dactylographié et relié au texte par une
accolade : « notion duale "coadmissible" : compatible avec Hom ».
« "admissible" ou » est dactylographié dans l'interligne ; une première
frappe, recouverte de x, précède $(P) \otimes_{o} F(X)$}
\marginal{[$P \otimes_{o} X \simeq \mathrm{Hom}_{o}(\check{P}, X)$]}
\note{ajout à la main au-dessus de l'homomorphisme naturel}
Si $C$ est une catégorie de Deuring, alors il suffit de poser cette
condition pour un objet fixé $X$ de $C$, et pour $o = \mathrm{End}_{C}(X)$,
$P$ un module $\in \mathbb{P}(o/A) =$ catégorie des $o$-modules à droite
[\uncertain{équiv.}] qui sont localement sur $S = \operatorname{Spec}(A)$
libres de rang~$1$. Dans ce cas, la catégorie des foncteurs $F$ en question
est équivalente, par
\[
  F \longmapsto F(X) : \mathrm{Hom}_{A\text{-adm}}(C, C')
  \hookrightarrow \mathrm{Mod}(o, C') ,
\]
à la catégorie des objets [$X'$] de $C'$ avec opération [$A$-linéaire] de
$o$, tels que $P \otimes_{o} X'$ existe dans $C'$ pour tout
$P \in \operatorname{Ob} \mathbb{P}(o/A)$. En tous cas, la catégorie des
foncteurs admissibles de $C$ dans $C'$ est $A$-linéaire.
\note{Hom est souligné deux fois, et $\mathbb{P}$ souligné, à la machine ;
un premier $\mathbb{P}(o/A)$ est frappé par-dessus après « un module
$\in$ »}

\marginal{il suffit que $C$ admette un $X$ « générateur strict » i.e.\ tout
objet de $C$ \uncertain{est} de la forme $P \otimes_{o} X$, avec $P$
projectif de type fini, alors $P \longmapsto P \otimes_{o} X$ (\ill{}
$\mathbb{P}(o)X$) \ill{} pour $P \otimes_{o} X$ existe)}
\note{note à la main dans la marge de gauche, à l'encre brune, écrite en
oblique à hauteur de « est un isomorphisme. Si $C$ … » ; elle annonce le
« générateur strict » du feuillet 62 (lot 4)}

Soit
\[
  C(1) = \mathrm{Hom}_{A\text{-adm}}(C^{\circ}, \mathrm{Mod}(A)) ,
\]
on a alors un foncteur pleinement fidèle canonique
\[
  C \xrightarrow{\ i\ } C(1)
\]
associant à tout objet $X$ de $C$ le foncteur
$Y \longmapsto \mathrm{Hom}_{C}(Y, X)$ qu'il représente.
\note{l'indice « A-adm » de ce Hom est frappé par-dessus un premier indice
recouvert ; dans la marge, à la main : « coadm »}

[Si $X$ est un objet de $C$ tel que tout objet $Y$ de $C$ s'écrit sous la
forme $P \otimes_{o} X$, avec $P$ un $o$-module à droite projectif de type
fini ($o = \mathrm{End}_{C}(X)$).]
\note{ces trois dernières lignes sont à la main, au bas du feuillet
dactylographié}

\page{59}
\note{le feuillet est entièrement à la main, sur le papier du tapuscrit ; il
continue la phrase manuscrite qui clôt le feuillet 58}
Alors pour tout objet $Y$ de $X$, \struck{$i(X)(Y)$} posant $F = i(X)$,
$F(Y)$ est un module [à droite] projectif de type fini sous
$o'^{\circ}$, où $o' = \mathrm{End}(Y)$. Soit \struck{$C_{sp}$}, dans le
cas général, $C_{sp}(1)$ la catégorie de ces foncteurs [(dits
« \uncertain{spéciaux} »)]. On trouve donc ici une inclusion (sans réserver
d'existence des gén.\ strict)
\[
  C \xrightarrow{\ i\ } C_{sp}(1)
\]
et le choix d'un \struck{générateur} [strict] $X$ de $C$ permet
d'identifier $C_{sp}(1)$ à la catégorie des $o$-modules à droite
projectifs de type fini. \struck{Et} On l'appelle la
\struck{\ill{}} [\uncertain{saturée-projective}] \struck{\ill{}}
[\uncertain{simple}] (ou la saturée) de $C$.
\note{« de $X$ » à la première ligne est écrit ainsi, pour « de $C$ »
peut-être. Dans la marge de gauche, à hauteur de la définition de
$C_{sp}(1)$, une note de cinq lignes biffée de traits obliques, \ill{}}
\marginal{il faut \uncertain{être prudent} : \ill{} partout : loc.\ libre
de rang fini sur $\operatorname{Spec} A = S$}
\note{note à l'encre brune dans la marge de gauche, en bas, écrite en
oblique ; la moitié inférieure du feuillet est blanche}

\page{60}
(1) Soit $C$ une catégorie $A$-linéaire. Soit \struck{\ill{}} $X$ un objet
de $C$, avec op.\ d'une $A$-algèbre $o$, et soit $P$ un $o$-module
[à droite] [projectif de type fini], si $P \otimes_{o} X$ existe
[\uncertain{(que $P$ soit)}] dans $C$, on a un hom.\ $o$-linéaire
\uncertain{naturel}
\[
  \check{P} \longrightarrow \mathrm{Hom}(P \otimes_{o} X, X) ,
\]
d'où un hom.\ de [contrav.]\ foncteurs \struck{\ill{}} sur $C$ [i.e.\ sur
$C^{\circ}$]
\[
  P \otimes_{o} X \longrightarrow \mathrm{Hom}_{o}(\check{P}, X)
\]
\note{à droite de la formule, en petit : « $P \otimes_{o} X$ » suivi d'un
trait qui renvoie à la ligne suivante}
On dit que le produit tensoriel externe est strict si la flèche précédente
est un isom. On dit que \struck{l'objet $P$} « les produits tensoriels
externes de $C$ sont raisonnables » si la condition précédente est
satisfaite chaque fois que $P$ est \struck{projectif} [projectif] de type
fini ; [\uncertain{Ceci} si $C$ est additive et karoubienne,]
\struck{\ill{} suffit pour \ill{}} dans cette condition est satisfaite, et
la $P \otimes_{o} X$ existe pour tout $P$ projectif de t.f.\ sur $o$.
[Il \struck{condition} suffit de l'exiger pour $o = \mathrm{End}(X)$.]
Alors le foncteur $P \longmapsto P \otimes_{o} X$ sur la sous-catégorie des
$P$ proj.\ de t.f.\ tels que [$P \otimes_{o} X$ existe, est pleinement
fidèle.]
\marginal{(condition automatique)}
\note{la note marginale, à l'encre brune, est placée à hauteur de « de type
fini ; » ; la fin de la phrase, « $P \otimes_{o} X$ existe, est pleinement
fidèle », est écrite dans la marge de gauche et renvoyée ici par une
flèche}

(2) $C$ est une catégorie \struck{de Deuring (resp.\ de \ill{})}
\struck{\ill{}} \emph{jolie} si \struck{(resp.\ \ill{} catégorial) si les
$P \otimes$} « les produits tensoriels y sont raisonnables » et si pour
deux objets $\{X, X'\}$, $\mathrm{Hom}(X, X')$ comme \struck{bimodule}
module sous $o$ [$o'$] \struck{$o$ et comme module sous $o'$} est projectif
de t.f.\ [et admet $o'$ comme commutant]. Si $C$ admet un gén.\ strict $X$,
(tout objet
\note{la phrase se poursuit au feuillet 62 (lot 4)}

\end{document}
